%! TeX program = lualatex
\documentclass{../../Styles/AI-Research-STY}
\addbibresource{../../References/references.bib}

\setdefaultlanguage{spanish}

\usepackage[
    pdftitle={AI Jugando Las Traes},
    pdfauthor={Roberto André Mossi Milla},
    pdfsubject={Inteligencia artificial multiagente mediante aprendizaje evolutivo y por refuerzo},
    pdfcreator={LaTeX Y Neovim}
    ]{hyperref}

\begin{document}
\headerBlock{
    Published in the Proceedings of the ASEE Annual Conference\textsuperscript{1}\par
    Presentado en la Universidad Estatal de Pensilvania (Sigma Xi)\textsuperscript{2}
}
\begin{center}
	\researchTitle{AI Jugando Las Traes}

	\authorInfo{Roberto André Mossi Milla}{
		Estudiante de Ciencias de la Computación, Universidad de Gannon
	}
	\authorInfo{Dr. Ramakrishnan Sundaram}{
		Profesor de Ingeniería Eléctrica y Cibernética, Universidad de Gannon
	}

	\vspace{4cm}
\end{center}
\begin{Form}

    \renewcommand{\AbstractTitle}{Resumen}
	\abstractSection{
        Este artículo desarrolla un AI que juega a las traes a través de la competencia entre diferentes sujetos y el aspecto de
		supervivencia de un juego de \textit{las traes}. Mediante el uso de genomas, se crean entidades sofisticadas que se
		coordinan entre sí para el beneficio de su supervivencia. A través de múltiples pruebas utilizando aprendizaje
		reforzado con genomas como base, se buscan los genes de supervivencia más aptos y estrategias de auto-desarrollo.
		Posteriormente, se introducen sus capacidades de razonamiento mediante el uso de entornos más complejos generados
		aleatoriamente, lo que conduce a habilidades más similares a las humanas. Esta investigación ya ha sido realizada
		previamente en otros juegos como Open-AI hide \& seek, pero para este estudio se añadieron comportamientos más
		humanos y se utilizaron genomas para observar su ingenio desde una perspectiva biológica. Esta investigación
		apoya el desarrollo del currículo al introducir a estudiantes del Departamento de Ciencias de la Computación (CS)
		al desarrollo de inteligencia artificial, y cómo otras áreas como la biología pueden aportar un desarrollo
		adicional a la tecnología que desarrollan. Se compara cómo esto se relaciona con la estructura biológica de los
		organismos en su toma de decisiones basada en tareas, inicialización aleatoria y restricciones. Con base en su
		comportamiento, se propondrán métodos de ajuste fino para mejorar la investigación hacia una replicación más
		cercana del comportamiento humano mediante el uso de genomas.
	}

    \footnotetext[1]{Publicación de ASEE disponible en \citeurl{AITagSandboxPaper}}

    \footnotetext[2]{Presentado en la Exposición de Investigación Sigma Xi de la Universidad Estatal de Pensilvania
    (\url{https://sites.psu.edu/behrendsigmaxi2025/2025/03/11/tag-ai-sandbox/})}

    \ClearShipoutPicture

	\midSection{Introducción}
	\midSectionContent{
		Utilizando genomas, el comportamiento de los sujetos será determinado a través de sus genes.
		Los genes se determinan al inicio mediante un proceso aleatorio, con el fin de otorgar a las entidades un punto
		de partida para aprender y adaptarse a las tareas de supervivencia que deben completar. A partir de ahí, se les
		asigna una tarea; en este caso, se establecen aleatoriamente dos grupos, los cuales se dividen en corredores y
		buscadores. La tarea de estos grupos consiste en que los buscadores deben atrapar a los corredores, y los
		corredores deben huir de ellos en un número limitado de intentos. Los buscadores que logran atrapar a su objetivo
		sobreviven, mientras que aquellos que no lo hacen mueren inmediatamente. El equipo contrario aplica el mismo
		principio, pero con el objetivo de evitar ser atrapado.

		Para el desarrollo completo de este proyecto, se utilizó Unreal Engine. La razón detrás de usar Unreal Engine
		es que las interacciones de comportamiento similares a las humanas requieren física. Existen muchas opciones,
		pero Unreal Engine es uno de los motores más populares del mercado, aunque también es conocido por su mayor
		complejidad debido a que su lenguaje principal es C++. Se tomó esto como un desafío con la esperanza de obtener
		el mayor rendimiento posible. Para lograrlo, se cuenta con la asistencia de varias librerías. Una de ellas es
		PyTorch, utilizada para el desarrollo de inteligencia artificial. La otra es CUDA, cuya funcionalidad consiste
		en redirigir el procesamiento gráfico del proyecto a la tarjeta gráfica, permitiendo que el resto del sistema
		se encargue del procesamiento del comportamiento.
	}

	\midSection{Trabajos Relacionados}
	\midSectionContent{
		Durante un tiempo, han existido múltiples intentos de crear inteligencia artificial enfocada en el estudio de
		organismos vivos. Para este caso específico, existe inspiración en dos proyectos concretos. El primer proyecto
		es \citetitle{MultiAgentHideandSeek} \parencite{MultiAgentHideandSeek}, desarrollado por Open-AI. Este proyecto
		brindó la inspiración para crear un entorno basado en un juego con el fin de estudiar el ingenio de la IA en
		situaciones de resolución de problemas. El segundo proyecto es el video
		\citetitle{ProgrammedSomeCreatures} \parencite{ProgrammedSomeCreatures}. Este proyecto tiene como idea principal
		la gestión del comportamiento de los organismos mediante el uso de genomas bajo las restricciones del entorno.
		Esto establece una conexión entre el mundo de la informática y la biología.
	}

    \midSection{Requisitos}
    \midSectionContent{
        \tableBlock{|l|l|l|l|l|}{
            Tipo     & Versión & Nombre & Descripción                      & Razón
        }{
            Lenguaje & C++20   & C++    & Lenguaje OOP de bajo nivel       & Es ampliamente utilizado en      \\
                     &         &        & usado para Unreal Engine         &  programación de IA              \\
                     &         &        &                                  & y por su velocidad.              \\
                     \hline
            Lenguaje & 3.31.6  & CMake  & Lenguaje utilizado para          & Al usar múltiples repositorios,  \\
                     &         &        & automatización entre sistemas    & se emplea CMake para automatizar \\
                     &         &        & operativos                       & la configuración y asegurar      \\
                     &         &        &                                  & compatibilidad multiplataforma   \\
                     \hline
            Librería & 12.0.0  & Cuda   & Modelo de programación para GPU  & Para procesamiento gráfico       \\
                     &         &        &                                  & intensivo, se usa esta librería  \\
                     &         &        &                                  & para acelerar el cálculo por GPU \\
                     \hline
            Gráficos & 5.4.4   & Unreal & Creación gráfica y visualización & Capaz de manejar computación     \\
                     &         & Engine & de entidades                     & compleja y alto rendimiento
                 }
             }

	\midSection{Generación del Mundo}
	\midSectionContent{
		En este caso era evidente que las entidades requerían un mundo en el cual pudieran desplazarse.
		Inicialmente se consideró crear un mundo estático genérico donde pudieran moverse libremente y aprender el
		juego por sí mismas. Sin embargo, surgió la posibilidad de que eventualmente se volvieran muy buenas jugando,
		pero sin espacio adicional para crecer al estar limitadas al mismo entorno. Para resolver esto, se determinó
		que sería necesario generar nuevos mundos para que pudieran aplicar lo aprendido en nuevos escenarios y así
		mejorar su toma de decisiones. Hacer esto manualmente después de cada ciclo implicaría un costo de tiempo
		significativo, por lo que se decidió implementar un generador de mundos\footnote{El código puede encontrarse en https://github.com/AndreM222/Procedural-Generator}. Esta decisión ya había sido aplicada en proyectos a largo plazo como Hide \& Seek de Open-AI, aunque en proyectos pequeños suele modelarse el mundo manualmente. Para evitar redundancia, un generador aleatorio resulta beneficioso.

		Para crear este mundo, primero se planificó el algoritmo a utilizar. Se revisaron diversos recursos sobre
		generadores en Unreal Engine y fuera de él. Muchos utilizaban el algoritmo A-Star, basado en el algoritmo de
		Dijkstra, el cual considera vecinos, rutas más cortas y distancia al nodo final \parencite{AStarAlgorithm}. Esto
		es útil cuando existe un punto inicial y final definido \parencite{IRewroteMyDungeonGenerator}.

		En su lugar, se optó por un enfoque similar a un algoritmo de búsqueda en anchura \parencite{BreadthFirstSearch},
		sin un punto final definido, sino con un número total de salas a generar. Dado que el objetivo del juego no es
		llegar a un destino específico, sino huir o perseguir, se eligió una solución más aleatoria.

		El proceso funciona de la siguiente manera: se define un punto inicial llamado origen, el cual genera una sala
		aleatoria. Las salas contienen flechas de salida que indican posibles conexiones. El generador elige aleatoriamente
		una salida y genera la siguiente sala conectada, repitiendo el proceso hasta alcanzar el número deseado de salas
		\parencite{HowToCreateEpicProceduralDungeons}.

		\imageBlock{../../IMG/ExitsPreview.png}{Configuración de Salas}
		\imageBlock[0.25]{../../IMG/ProceduralPreviewPic.png}{Vista Previa del Generador}

		\midSubSection{Detección de Colisiones}
		\midSubSectionContent{
			¿Pero qué ocurrirá si ya no hay más salidas disponibles? Teniendo esa pregunta en consideración, creé dos funciones
			adicionales: reboot y delete room. La función reboot lo que hace es ejecutar la eliminación de cada habitación que
			haya sido creada; una vez que todas han sido eliminadas, ejecuta una función de borrado para el generador actual,
			la cual espera a que se ejecuten todas las instrucciones restantes antes de eliminarse a sí misma. Después de eso,
			se ejecuta la última instrucción, que genera un nuevo generador de habitaciones e inicia todo el proceso nuevamente,
			generando aleatoriamente un nuevo mundo con los requisitos dados.

			Pero, ¿qué sucede con la colisión entre diferentes habitaciones?, ¿se superpondrán? Para esta situación creé
			varias funciones que cubren colisiones en distintos escenarios. Estas situaciones pueden separarse en dos
			categorías diferentes.

            \begin{bulletBlockContent}
                \bulletBlockItem{Colisión General} Para la colisión general, todo lo que se requiere es obtener el espacio donde
                se intenta generar una habitación y comprobar si existe algo en ese punto. Para ello, la función genera una
                caja de trazado invisible que, si colisiona con cualquier objeto, indica que se debe buscar una nueva ubicación.
            \end{bulletBlockContent}
		}
	}

    \begin{codeBlock}{cpp}[Comprobación de barrido de colisión]
// Setup Box Dimensions for trace
FVector BoxDimensions = FVector(50, 50, 50);

FVector traceLocation = exitLocation.GetLocation();

// Trace box to check if collision available
FCollisionShape box = FCollisionShape::MakeBox(BoxDimensions);
exitsList[spaceIndex]->GetWorld()->SweepSingleByChannel(
    Hit,
    traceLocation,
    traceLocation,
    FQuat::Identity,
    ECC_Visibility,
    box,
    QueryParams
);

// Return true if room type has been found connected
if (Cast<AMaster_Room>(Hit.GetActor())) return true;
    \end{codeBlock}

	\midSectionContent{
        \begin{bulletBlockContent}
            \bulletBlockItem{Generación por Dimensiones} En base a las dimensiones de la habitación que se intenta generar,
            es necesario verificar si existe espacio suficiente para que dicho punto pueda ser ocupado. Sin embargo, obtener
            las dimensiones de una habitación de antemano no es posible con la versión actual de Unreal Engine. Para resolver
            esto, implementé una solución improvisada: generar brevemente la habitación deseada para obtener sus dimensiones,
            almacenarlas en una variable de transformación y luego destruir la habitación. Todo este proceso ocurre tan rápido
            que la habitación generada no llega a ser visible. A partir de ahí, junto con la colisión general, se pasan los
            nuevos valores calculados y se reemplazan las dimensiones de la caja de trazado para realizar una prueba de
            colisión más precisa.
        \end{bulletBlockContent}
	}

    \begin{codeBlock}{cpp}[Generador de dimensiones]
// Make temporary room
AActor* roomModel = GetWorld()->SpawnActor<AMaster_Room>(
    RoomToSpawn[roomIndex],
    exitLocation.GetLocation(),
    FRotator::ZeroRotator,
    SpawnParams
);

// Setup Box Dimensions for trace
FVector BoxDimensions = FVector(roomModel->GetComponentsBoundingBox().GetExtent() - 4);
BoxDimensions.Z += 3;
roomModel->Destroy(); // Destroy tmp room after setup

FVector traceLocation = exitLocation.GetLocation() + (
    exitLocation.GetLocation() - actorLocation.GetLocation());
traceLocation.Z += BoxDimensions.Z; // Center Box Trace

// Trace box to check if collision available
FCollisionShape box = FCollisionShape::MakeBox(BoxDimensions);
    \end{codeBlock}

	\midSubSection{Encapsulación del Mundo}
	\midSubSectionContent{
		En el mundo que estamos creando, es necesario agregar ciertas restricciones para las entidades. Estas restricciones
		ayudan a evitar que las entidades salgan de los límites del mundo y caigan al vacío. El enfoque principal de este
		proyecto no es que aprendan dónde pueden caerse, sino cómo reaccionar ante la interacción entre diversas entidades.
		Para estas prevenciones, decidí encapsular con paredes las salidas que quedaron vacías al final de la generación.

		Al cerrar las paredes, es importante evitar que se generen en lugares que conectan con otras habitaciones. Si no se
		tomara esta precaución, no habría diferencia entre un mundo y un laberinto. Para aclarar esta distinción: un laberinto
		es generalmente restrictivo; el movimiento está limitado por caminos delimitados por paredes, comúnmente definidos
		como corredores. En cambio, un mundo, en este contexto, es menos restrictivo y está compuesto tanto por corredores
		como por espacios abiertos.

		Para evitar que las paredes se generen en todos los lugares posibles, se crean pequeñas cajas de trazado. Estas cajas
		se utilizan únicamente para detectar si existe alguna habitación conectada al punto donde se desea colocar la pared.
		Si se detecta una habitación conectada, dicha salida se elimina de la lista de salidas abiertas y se continúa con la
		siguiente salida a cerrar. Si la habitación deseada no puede colocarse en ninguna de las salidas disponibles, se
		elimina de la lista de habitaciones posibles y se continúa con la siguiente hasta que no queden más salidas ni
		habitaciones disponibles.

		\imageBlock[0.4]{../../IMG/MasterRoomChildrenPreview.png}{Creación de Habitaciones}
	}

	\midSubSection{Semilla}
	\midSubSectionContent{
		En cierto punto me di cuenta de que podría ser útil mantener a las entidades en una habitación específica durante un
		período de tiempo. Para ello decidí agregar una opción de semilla\footnote{Más información disponible en https://github.com/AndreM222/Procedural-Generator/blob/main/docs/setup.md}.
		Cuando se utiliza una semilla, dependiendo del valor, se generará una habitación que permanecerá igual sin importar
		cuántas veces se recargue. Sin embargo, el uso de una semilla introduce la posibilidad de que no se logre generar el
		número total de habitaciones deseado. Este fenómeno puede ocurrir cuando las habitaciones se generan siguiendo un
		camino específico que eventualmente termina cerrándose sobre sí mismo, sin dejar salidas disponibles. Para este
		caso, indiqué al generador que omitiera el ciclo de creación de nuevos generadores hasta que uno cumpliera con los
		requisitos establecidos. Esto se hizo para evitar la ocurrencia de un bucle infinito.

		\imageBlock{../../IMG/SpawnPointsPreview.png}{Generador de Props}

		\imageBlock{../../IMG/PropsPreview.png}{Configuración de Props}
	}

	\midSubSection{Props y Generadores}
	\midSubSectionContent{
		Para añadir generadores de entidades y obstáculos, implementé una función de generación que crea actores de manera
		aleatoria con capacidades de filtrado, distribuyéndolos a lo largo de las habitaciones. Esta función busca las flechas
		de generación dentro de la carpeta de generadores.
	}

	\midSection{Configuración del Personaje}
	\midSectionContent{
		Para el modelo 3D y las animaciones estoy utilizando los recursos de un paquete llamado \citetitle{AdvanceLocomotionSystemV4} \parencite{AdvanceLocomotionSystemV4}. Este
		paquete contiene una gran cantidad de animaciones útiles y algunos modelos simples de los cuales pude beneficiarme
		para enfocarme en el desarrollo de las acciones y el razonamiento de los personajes.

		Este paquete también incluye una configuración para el movimiento, pero fue realizada completamente en blueprints,
		lo cual no es fácilmente accesible desde C++. Para esta situación desarrollé mi propia lógica con respecto a la
		interacción que las entidades\footnote{El código puede encontrarse en https://github.com/AndreM222/AI-Entities} realizan
		en el espacio 3D. Los comportamientos incluidos en el proyecto fueron realizados completamente utilizando blueprints.
		Esta es una de las ventajas de Unreal Engine, pero surge un problema. Actualmente no es posible interactuar desde C++
		hacia los blueprints, aunque sí puede hacerse de forma inversa. Por lo tanto, desarrollé mi propia versión de
		la locomoción para lograr un movimiento realista, utilizando las animaciones del paquete.

		\imageBlock[0.6]{../../IMG/Teams-Picture-Preview.png}{Configuración del Personaje}

		\midSubSection{Estados}
		\midSubSectionContent{
			Los estados son las diferentes posturas en las que las entidades pueden cambiar dependiendo de su interacción,
			con el objetivo de aplicar interacciones más humanas en el mundo real. Las interacciones podrían haberse realizado
			sin animaciones, pero si se desean considerar los retrasos en la reacción y una mejor facilidad de visualización,
			decidí utilizar las animaciones del paquete de recursos para hacerlo posible. A continuación se muestran los
			estados disponibles creados.

			\begin{bulletBlockContent}
				\bulletBlockItem{Usado} Actualmente se utiliza para algún tipo de interacción.
				\bulletBlockItem{Planificado} Está listo para ser utilizado, pero aún no se ha combinado con la lógica para activarlo.
			\end{bulletBlockContent}

            \tableBlock{|l|l|l|}{
                Postura    & Descripción                                                  & Estado
            }{
                Default    & Es una postura que invoca neutralidad y paz                  & Planificado \\
                \hline
                Masculine  & Está listo para combate o para mostrar poder                 & Usado       \\
                \hline
                Feminine   & Es una postura más elegante y noble                          & Planificado \\
                \hline
                Injured    & Muestra la baja cantidad de salud de las entidades           & Usado       \\
                \hline
                HandsTied  & Muestra cuando existe una restricción en las entidades       & Planificado \\
                \hline
                Rifle      & Sosteniendo un rifle                                         & Planificado \\
                \hline
                Pistol 1H  & Sosteniendo una pistola con una mano                         & Planificado \\
                \hline
                Pistol 2H  & Sosteniendo una pistola con ambas manos                      & Planificado \\
                \hline
                Bow        & Sosteniendo un arco                                          & Planificado \\
                \hline
                Torch      & Sosteniendo una antorcha para iluminación                    & Planificado \\
                \hline
                Binoculars & Sosteniendo binoculares para ver a larga distancia           & Planificado \\
                \hline
                Box        & Cargando cajas                                               & Planificado \\
                \hline
                Barrel     & Cargando un barril                                           & Planificado \\
                \hline
                Reach      & Postura para alcanzar y tocar el objetivo a corta distancia  & Planificado
            }
        }
    }

	\midSection{Acciones}
	\midSectionContent{
		Las acciones son las diferentes capacidades que las entidades tienen para interactuar con el mundo. Esto añade
		mayor complejidad y variedad al razonamiento de las entidades.

        \tableBlock{|l|l|l|}{
            Acción & Descripción                    & Estado
        }{
			Run    & Movimiento a alta velocidad    & Usado       \\
			\hline
			Sprint & Movimiento a velocidad media   & Usado       \\
			\hline
			Walk   & Movimiento a baja velocidad    & Usado       \\
			\hline
			Crouch & Agacharse en espacios pequeños & Usado       \\
			\hline
			Climb  & Escalar paredes                & Usado       \\
			\hline
			Grab   & Agarrar un objeto              & Planificado \\
			\hline
			Drop   & Soltar el objeto sostenido     & Planificado \\
			\hline
			Jump   & Saltar                         & Usado       \\
			\hline
			Touch  & Tocar el objetivo              & Usado
		}

		\imageBlock[0.4]{../../IMG/ActionPreview.png}{Vista Previa de Acciones}
	}

	\midSection{Toma de Decisiones a Través de Genomas}
	\midSectionContent{
		He comenzado el desarrollo de la inteligencia de los personajes utilizando genomas, pero aún faltan una variedad
		de funciones para que puedan tomar decisiones racionales a través de generaciones de datos proporcionados por
		los padres.

		Esta aplicación tuvo éxito en tareas simples, como que las entidades tengan la necesidad de dirigirse a áreas
		específicas del mapa. Sin embargo, el juego de atrapar (las traes) es una tarea más compleja que requerirá más datos y
		capacidades por parte de las entidades.

		\midSubSection{Algoritmo Genético}
		\midSubSectionContent{
			Para más detalles sobre cómo funciona esto, se utiliza un algoritmo genético \parencite{GeneticAlgorithm}, lo cual
			implica el uso de la evolución. En cuanto a la toma de decisiones, podríamos permitir que una entidad intente
			cablear su propio cerebro mediante algo llamado mutación. Sin embargo, habrá un límite en cuanto a cuánto puede
			crecer o encontrar mejores formas de resolver un problema. Como todos sabemos, cuando se trata del desarrollo
			personal existen factores internos y externos que nos ayudan a crecer. Los factores internos se entienden como
			las experiencias que hemos vivido y que utilizamos para tomar decisiones, mientras que los factores externos
			son aquellos que nos enseñan a partir de su propia experiencia. Este factor no solo nos hace crecer, como lo hace
			el interno, sino que nos permite crecer de la manera más eficiente.

			\imageBlock[0.15]{../../IMG/GeneticAlgorithm.png}{Algoritmo Genético}
		}

		\midSubSection{Algoritmo NEAT}
		\midSubSectionContent{
			El algoritmo NEAT \parencite{NeatAlgorithm} utiliza el algoritmo genético, pero pone en práctica la idea de los
			genomas para desarrollar sus cerebros. Al contar con sentidos, acciones y neuronas internas, estas neuronas se
			conectan como topologías aumentadas, lo que significa que su cerebro evoluciona según sus necesidades. Sin
			embargo, para mantener cierto control, se establece un límite máximo de neuronas.

			El algoritmo toma dos padres y los combina utilizando feromonas determinadas por su similitud, la cual en este
			caso se define como qué tan bien resuelven su tarea. Todas estas neuronas pueden combinarse entre sí y consigo
			mismas. Para realizar una limpieza, se implementó una función que detecta si alguna conexión no cumple ninguna
			función y la elimina.

			Si una entidad tiene demasiados sentidos, puede verse abrumada y no ejecutar correctamente sus acciones, pero si
			tiene muy pocos, no encontrará la mejor forma de resolver un problema.

			\imageBlock[0.15]{../../IMG/NeatAlgorithm.png}{Algoritmo NEAT}
		}

		\midSubSection{Senses}
		\midSubSectionContent{
			Para estas entidades, comencé a desarrollar sentidos similares a los de los humanos. Por sentidos me refiero al
			sentido del oído, el sentido de la vista, etc. Sin embargo, para este proyecto limitaremos estos sentidos a los
			siguientes.

            \tableBlock{|l|l|l|}{
                Tipo de Sentido & Descripción                                                                  & Estado
            }{
				Espacio         & Este sentido permitirá a las entidades saber dónde se encuentran.            & En Uso        \\
				\hline
				Tacto           & Se utilizará para obtener datos de sus interacciones y decidir si            & En Desarrollo \\
				                & les agrada o desagrada.                                                      &               \\
				\hline
				Vista           & Permite a las entidades saber qué tienen frente a ellas y decidir qué hacer. & En Uso
			}

			Pero cada entrada tiene subdivisiones que crean complejidad adicional. Por ejemplo, para obtener la vista es
			necesario calcular la distancia a los actores, la distancia a las paredes, la cantidad de población, entre otros.
			Para ello, se utilizan diversos trazados con el fin de calcular distancias, cantidades, qué hay al frente, etc.

			\imageBlock[0.35]{../../IMG/TracingSense.png}{Vista Previa de los Sentidos}

			El comportamiento inicial solo contaba con correr, saltar y escalar como acciones. En cuanto a los sentidos,
			se detectaba la densidad de población y la distancia, así como una configuración mínima de la distancia a los
			límites. Los personajes comenzaron únicamente a moverse en círculos y no realizaban nada en particular. Por
			ello, se añadieron más sentidos y acciones como la detección de bordes, la ubicación actual basada en los
			límites y valores normalizados. Con esto se incrementó su proceso de razonamiento y comenzaron a navegar más;
			algunos no hacían nada, y otros saltaban si se detectaba una pared cercana o, por el contrario, si no la había.
			Resultó fascinante observar cómo comenzaron a desarrollarse más a medida que podían percibir y realizar más
			acciones.
		}
	}

	\midSection{Todo en Uno}
	\midSectionContent{
		Dado que la escala de este proyecto es bastante grande y que además muchos recursos pueden utilizarse para pruebas
		que no serán necesarios en el producto final, separé todos los componentes en distintos repositorios. Sin
		embargo, esto implica que en algún punto debe realizarse una fusión. Unreal Engine no es tan sencillo de manejar
		cuando se trata de fusionar distintos proyectos en uno solo, a diferencia de otras herramientas, por lo que fue
		necesario crear un sistema de fusión automatizado utilizando un archivo de cmake
		\footnote{Para el proyecto fusionado, véase https://github.com/AndreM222/Tag-AI-Sandbox}.

		Este proyecto se encargará de la interfaz de usuario y de la comunicación entre todos los componentes creados sin
		afectar los scripts originales. Esto permite mantener la seguridad del código al aislarlo en diferentes
		repositorios, evitando así posibles rupturas.

		\midSubSection{Por qué CMake}
		\midSubSectionContent{
			Inicialmente pensé en crear un archivo bash para ejecutar la fusión entre los proyectos, pero esto habría
			limitado la configuración a sistemas operativos específicos. CMake, aunque requiere una instalación, permite
			que el proceso de ejecución sea multiplataforma, lo que lo convierte en una opción ideal.
		}

		\midSubSection{Cómo Añadir}
		\midSubSectionContent{
			Para facilitar la configuración, decidí utilizar un archivo json para listar los proyectos que se desean
			fusionar, así como una opción de filtrado. Este sistema de fusión se encarga de eliminar archivos no importantes,
			como los de compilación y el script principal de los proyectos, y reemplaza todas las APIs por las del proyecto
			actual.
		}
	}

	\midSection{Evaluación}
	\midSectionContent{
		\midSubSection{Trayectoria}
		\midSubSectionContent{
			Las entidades parecen terminar siguiendo un patrón de desplazarse hacia direcciones específicas hasta que
			logran tocar su objetivo o evitarlo. Sin embargo, se observa una mejora constante con el tiempo en su
			comportamiento y adaptación a diferentes entornos.

			Intenté ejecutarlo bajo los mismos requisitos que \citetitle{ProgrammedSomeCreatures}
			\parencite{ProgrammedSomeCreatures} para comprobar qué ocurriría con una misión más básica, como dirigirse a
			áreas específicas del mapa. Al ejecutarlo, terminaron reaccionando de la misma manera. Esto puede deberse a que
			un juego de atrapar requiere sentidos como la vista, el tacto y el sentido del espacio para tomar decisiones más
			racionales.
		}

		\midSubSection{Obstáculos}
		\midSubSectionContent{
			Dado que aún no se han realizado muchas pruebas, por ahora las entidades parecen caminar hacia las paredes en
			ciertos ángulos hasta que logran cruzar al otro lado si es necesario.

			Se requerirán más pruebas para observar cómo mejora su comportamiento con el tiempo.
		}
	}

	\midSection{Qué Sigue}
	\midSectionContent{
		Continuaré desarrollando la estructura para la toma de decisiones y trasladaré el proyecto a la supercomputadora
		para comenzar a obtener datos mientras continúo implementando nuevas funcionalidades. También necesito
		desarrollar una interfaz de usuario que me permita obtener datos de los eventos, así como tomar capturas de
		pantalla en cada generación. Mientras las entidades se ejecutan en segundo plano, estaré trabajando en los
		comportamientos adicionales para implementarlos.

		Para mejorar la toma de decisiones de las entidades, comenzaré a incorporar más sentidos. Esto mejorará su
		capacidad de decisión al poder elegir ubicaciones basándose en dónde estuvieron, lo que tocaron o lo que vieron,
		así como en lo que consideran interesante, peligroso o irrelevante.

		El comportamiento de las entidades aún es primitivo y requiere mayor desarrollo para observar mejoras más
		significativas. Se necesitan iteraciones adicionales para analizar cómo interactúan con sus entornos y si logran
		aprender exitosamente a atrapar o escapar del objetivo.
	}

\end{Form}

\bibliographyBlock{Referencias}

\end{document}
